\subsection{Tensor Products and Universal Properties}

% This subsection integrates the bilinear/tensor-product core. Symmetric,
% exterior, and iterated tensor constructions are developed in the next file.

Tensor products are the device that turns bilinear data into linear data. The guiding rule is the same one already used for quotient spaces in Theorem \ref{thm:quotient-universal-property}: a universal object is recognised by the unique way every compatible map factors through it.

\begin{definition} \label{def:finite-coordinate-tensor-product-model} \hlt{Finite-Coordinate Tensor Product Model}\\
Let $\beta=(\mathbf{v}_1,\ldots,\mathbf{v}_m)$ and $\gamma=(\mathbf{w}_1,\ldots,\mathbf{w}_n)$ be ordered bases of $\mathcal{V}$ and $\mathcal{W}$.\\
Define the \hlt{coordinate tensor model} $\mathcal{V}\otimes_{\beta,\gamma}\mathcal{W}\defeq\operatorname{Mat}(m,n,\mathbb{F})$ by $\mathbf{v}_i\otimes_{\beta,\gamma}\mathbf{w}_j\defeq E_{i,j}$.\\
If $\mathbf{v}=\sum\limits_{i=1}^m x_i\mathbf{v}_i$ and $\mathbf{w}=\sum\limits_{j=1}^n y_j\mathbf{w}_j$, then $\mathbf{v}\otimes_{\beta,\gamma}\mathbf{w}\defeq \sum\limits_{i=1}^m\sum\limits_{j=1}^n x_iy_jE_{i,j}\in\operatorname{Mat}(m,n,\mathbb{F})$.
\end{definition}

\begin{proposition} \label{prop:tensor-product-matrix-basis-model} \hlt{Tensor Model Basis and Bilinear Evaluation}\\
In the coordinate tensor model associated to bases $\beta=(\mathbf{v}_1,\ldots,\mathbf{v}_m)$ and $\gamma=(\mathbf{w}_1,\ldots,\mathbf{w}_n)$:
\begin{enumerate}[label=(\roman*)]
\setlength{\itemsep}{0pt}
\item the matrices $E_{i,j}=\mathbf{v}_i\otimes_{\beta,\gamma}\mathbf{w}_j$ form a basis;
\item the map $(\mathbf{v},\mathbf{w})\mapsto\mathbf{v}\otimes_{\beta,\gamma}\mathbf{w}$ is bilinear;
\item every bilinear form $B:\mathcal{V}\times\mathcal{W}\to\mathbb{F}$ determines a unique linear functional $\Theta_B$ on $\operatorname{Mat}(m,n,\mathbb{F})$ such that\\
$\Theta_B(\mathbf{v}\otimes_{\beta,\gamma}\mathbf{w})=B(\mathbf{v},\mathbf{w})$ for all $\mathbf{v}\in\mathcal{V}$ and $\mathbf{w}\in\mathcal{W}$.
\end{enumerate}
\begin{proof}
\begin{enumerate}[label=(\roman*)]
\setlength{\itemsep}{0pt}
\item This is the standard entry-matrix basis of $\operatorname{Mat}(m,n,\mathbb{F})$ from Proposition \ref{prop:matrix-spaces-are-vector-spaces}.
\item The coordinates of $a\mathbf{v}_1+\mathbf{v}_2$ are the same linear combination of the coordinates of $\mathbf{v}_1$ and $\mathbf{v}_2$, and the displayed formula is linear in those coordinates. The same argument applies in the second variable.
\item Define $\Theta_B(E_{i,j})=B(\mathbf{v}_i,\mathbf{w}_j)$ on the basis from (i), and extend linearly by Proposition \ref{prop:linear-map-determined-by-basis}. If $\mathbf{v}=\sum_i x_i\mathbf{v}_i$ and $\mathbf{w}=\sum_j y_j\mathbf{w}_j$, then $\Theta_B(\mathbf{v}\otimes_{\beta,\gamma}\mathbf{w}) =\sum\limits_{i=1}^m\sum\limits_{j=1}^n x_i y_j B(\mathbf{v}_i,\mathbf{w}_j) =B(\mathbf{v},\mathbf{w})$, by bilinearity of $B$.\\
Uniqueness follows since the values on basis matrices $E_{i,j}$ are forced by $\mathbf{v}_i,\mathbf{w}_j$.
\end{enumerate}
\hlt{Structural Commentary}: The coordinate tensor model stores each pair of basis directions in its own matrix slot, so a bilinear form can be evaluated by a single linear functional on those slots.\\
\hlt{Key Insight}: Tensor coordinates turn products of coordinates into linear coordinates.
\end{proof}
\end{proposition}

\begin{figure}[H]
  \centering
  \begin{tikzpicture}[
  >=Stealth,
  every node/.style={font=\small},
  map/.style={->, gray!65!black, line width=0.75pt},
  slot/.style={draw=gray!45, fill=gray!4, minimum width=0.68cm, minimum height=0.52cm},
  active/.style={slot, fill=teal!10, draw=teal!65!black, line width=0.65pt},
  basis/.style={draw=gray!45, fill=gray!3, rounded corners=3pt, line width=0.55pt,
    minimum width=2.45cm, minimum height=0.78cm, align=center}
]
  \node[basis] (pair) at (-2.45,0.52) {$\mathbf{v}_i\otimes_{\beta,\gamma}\mathbf{w}_j$};

  \draw[map] (pair.east) -- (-0.60,0.52);

  \foreach \x in {0,1,2,3} {
    \foreach \y in {0,1,2} {
      \node[slot] at (\x*0.70,\y*0.54) {};
    }
  }
  \node[active] at (1.40,0.54) {$E_{i,j}$};

  \draw[gray!50] (-0.02,0.54) -- (2.10,0.54);
  \draw[gray!50] (1.40,-0.26) -- (1.40,1.34);
\end{tikzpicture}

  \caption{A basis pair occupies the $(i,j)$ matrix slot.}
  \label{fig:tensor-product-matrix-basis}
\end{figure}

\begin{definition} \label{def:bilinear-map} \hlt{Bilinear Maps}\\
Let $\mathcal{U},\mathcal{V},\mathcal{W}$ be $\mathbb{F}$-vector spaces and $f:\mathcal{U}\times\mathcal{V}\to\mathcal{W}$.
\begin{enumerate}[label=(\roman*)]
\setlength{\itemsep}{0pt}
\item $f$ is \hlt{bilinear} if $\forall r,s\in\mathbb{F}$, $\mathbf{u},\mathbf{u}'\in\mathcal{U}$, $\mathbf{v},\mathbf{v}'\in\mathcal{V}$, $f(r\mathbf{u}+s\mathbf{u}',\mathbf{v})=r f(\mathbf{u},\mathbf{v})+s f(\mathbf{u}',\mathbf{v})$\\
and $f(\mathbf{u},r\mathbf{v}+s\mathbf{v}')=r f(\mathbf{u},\mathbf{v})+s f(\mathbf{u},\mathbf{v}')$.
\item $\operatorname{Bil}_{\mathbb{F}}(\mathcal{U},\mathcal{V};\mathcal{W})\defeq\{f:\mathcal{U}\times\mathcal{V}\to\mathcal{W}:f\text{ bilinear}\}$.
\item If $\mathcal{W}=\mathbb{F}$, then $\operatorname{Bil}_{\mathbb{F}}(\mathcal{U},\mathcal{V};\mathbb{F})$ is the space of \hlt{bilinear forms} from Definition \ref{def:bilinear-form-and-matrix}.
\end{enumerate}
\end{definition}

\begin{example} \label{ex:basic-bilinear-maps} \hlt{Basic Bilinear Maps}
\begin{enumerate}[label=(\roman*)]
\setlength{\itemsep}{0pt}
\item If $A$ is an $\mathbb{F}$-algebra, then $\mu:A\times A\to A$, $\mu(a,b)\defeq ab$, is bilinear.
\item Evaluation $\operatorname{ev}:\LL(\mathcal{V},\mathcal{W})\times\mathcal{V}\to\mathcal{W}$, $\operatorname{ev}(T,\mathbf{v})\defeq T\mathbf{v}$, is bilinear.
\item If $\varphi\in\mathcal{V}'$ and $\psi\in\mathcal{W}'$, then $B_{\varphi,\psi}:\mathcal{V}\times\mathcal{W}\to\mathbb{F}$, $B_{\varphi,\psi}(\mathbf{v},\mathbf{w})\defeq\varphi(\mathbf{v})\psi(\mathbf{w})$, is bilinear.
\end{enumerate}
\end{example}

\begin{definition} \label{def:tensor-product-universal-property} \hlt{Tensor Product by Universal Property}\\
Let $\mathcal{U},\mathcal{V}$ be $\mathbb{F}$-vector spaces.
\begin{enumerate}[label=(\roman*)]
\setlength{\itemsep}{0pt}
\item A pair $(\mathcal{T},t)$ is a \hlt{tensor product} of $\mathcal{U}$ and $\mathcal{V}$ if $\mathcal{T}$ is an $\mathbb{F}$-vector space, $t\in\operatorname{Bil}_{\mathbb{F}}(\mathcal{U},\mathcal{V};\mathcal{T})$, and for every $\mathbb{F}$-vector space $\mathcal{W}$ and $f\in\operatorname{Bil}_{\mathbb{F}}(\mathcal{U},\mathcal{V};\mathcal{W})$, $\exists!\widetilde{f}\in\LL(\mathcal{T},\mathcal{W})$ such that $f=\widetilde{f}\circ t$.
\item The map $\widetilde{f}$ is the \hlt{mediating map} for $f$.
\item Write $\mathcal{T}=\mathcal{U}\otimes\mathcal{V}$ and $\mathbf{u}\otimes\mathbf{v}\defeq t(\mathbf{u},\mathbf{v})$; such tensors are \hlt{decomposable}.
\end{enumerate}
\end{definition}

\begin{figure}[H]
  \centering
  % Source note: Roman, Advanced Linear Algebra, Ch. 14, Fig. 14.4.
% Redrawn as a schematic universal-property diagram for these notes.
\begin{tikzpicture}[
  >=Stealth,
  every node/.style={font=\small},
  arrow/.style={->, line width=0.55pt},
  factor/.style={->, dashed, line width=0.55pt},
  note/.style={font=\scriptsize, text=gray!70!black}
]
  \node (uv) at (0,1.25) {$\mathcal{U}\times\mathcal{V}$};
  \node (t) at (3.6,1.25) {$\mathcal{U}\otimes\mathcal{V}$};
  \node (w) at (3.6,-0.8) {$\mathcal{W}$};

  \draw[arrow] (uv) -- node[above, note] {$t$ bilinear} (t);
  \draw[arrow] (uv) -- node[below left, note] {$f$ bilinear} (w);
  \draw[factor] (t) -- node[right, note] {$\widetilde{f}$ linear} (w);
\end{tikzpicture}

  \caption{A tensor product is universal for bilinear maps out of $\mathcal{U}\times\mathcal{V}$.}
  \label{fig:tensor-universal-property}
\end{figure}

\begin{theorem} \label{thm:tensor-product-quotient-construction} \hlt{Coordinate-Free Quotient Construction}\\
Let $\mathcal{U},\mathcal{V}$ be $\mathbb{F}$-vector spaces and let $\mathbb{F}^{(\mathcal{U}\times\mathcal{V})}$ have basis $\{[\mathbf{u},\mathbf{v}]:(\mathbf{u},\mathbf{v})\in\mathcal{U}\times\mathcal{V}\}$.\\
Let $\mathcal{R}$ be the span of all $r[\mathbf{u}_1,\mathbf{v}]+s[\mathbf{u}_2,\mathbf{v}]-[r\mathbf{u}_1+s\mathbf{u}_2,\mathbf{v}]$; and $r[\mathbf{u},\mathbf{v}_1]+s[\mathbf{u},\mathbf{v}_2]-[\mathbf{u},r\mathbf{v}_1+s\mathbf{v}_2]$, where $r,s\in\mathbb{F}$, $\mathbf{u},\mathbf{u}_1,\mathbf{u}_2\in\mathcal{U}$, and $\mathbf{v},\mathbf{v}_1,\mathbf{v}_2\in\mathcal{V}$. Then $\mathcal{U}\otimes\mathcal{V}\defeq \mathbb{F}^{(\mathcal{U}\times\mathcal{V})}/\mathcal{R}$, $t(\mathbf{u},\mathbf{v})=[\mathbf{u},\mathbf{v}]+\mathcal{R}$ is a tensor product of $\mathcal{U}$ and $\mathcal{V}$.
\begin{proof} Let $\pi:\mathbb{F}^{(\mathcal{U}\times\mathcal{V})}\to\mathcal{U}\otimes\mathcal{V}$ be the quotient map.\\
Since the generators of $\mathcal{R}$ vanish in the quotient, $t(\mathbf{u},\mathbf{v})=\pi[\mathbf{u},\mathbf{v}]$ is bilinear.\\
Let $f\in\operatorname{Bil}_{\mathbb{F}}(\mathcal{U},\mathcal{V};\mathcal{W})$ and define $\alpha:\mathbb{F}^{(\mathcal{U}\times\mathcal{V})}\to\mathcal{W}$ by $\alpha[\mathbf{u},\mathbf{v}]\defeq f(\mathbf{u},\mathbf{v})$.\\
Bilinearity gives $\mathcal{R}\subseteq\ker\alpha$, so Theorem \ref{thm:quotient-universal-property} gives a unique $\widetilde{f}\in\LL(\mathcal{U}\otimes\mathcal{V},\mathcal{W})$ with $\widetilde{f}\circ \pi=\alpha$. Hence $\widetilde{f}\circ t=f$.\\
If $L\circ t=f$, then $L$ and $\widetilde{f}$ agree on the spanning set $\{[\mathbf{u},\mathbf{v}]+\mathcal{R}\}$, so $L=\widetilde{f}$.\\
\hlt{Structural Commentary}: The quotient kills precisely the failure of the raw ordered-pair basis to be bilinear.\\
\hlt{Key Insight}: A tensor product is a free bilinear map, followed by no extra identifications.
\end{proof}
\end{theorem}

\begin{figure}[H]
  \centering
  % Source note: Roman, Advanced Linear Algebra, Ch. 14, Fig. 14.5.
% Redrawn to show the free-vector-space quotient construction.
\begin{tikzpicture}[
  >=Stealth,
  x=0.9cm,
  y=0.9cm,
  every node/.style={font=\small},
  arrow/.style={->, line width=0.65pt, black},
  aux/.style={->, dashed, line width=0.65pt, black},
  factor/.style={->, dashed, line width=0.7pt, black},
  note/.style={font=\scriptsize, text=black, fill=white, inner sep=0.8pt}
]
  \node (uv) at (0,0.45) {$\mathcal{U}\times\mathcal{V}$};
  \node (free) at (3.2,0.45) {$\mathbb{F}^{(\mathcal{U}\times\mathcal{V})}$};
  \node (quot) at (6.55,0.45) {$\mathcal{U}\otimes\mathcal{V}$};
  \node (w) at (3.2,-0.58) {$\mathcal{W}$};

  \draw[arrow] (uv) -- node[above, note] {$j$} (free);
  \draw[arrow] (free) -- node[above, note, pos=0.55] {$\pi$} (quot);
  \draw[arrow] (uv.north east) to[out=30,in=150]
    node[above, note] {$t=\pi j$} (quot.north west);
  \draw[arrow] (uv) -- node[below left, note] {$f$} (w);
  \draw[aux] (free) -- node[right, note, xshift=2pt] {$\alpha$} (w);
  \draw[factor] (quot) -- node[right, note, xshift=8pt] {$\widetilde{f}$} (w);
\end{tikzpicture}

  \caption{The free linear map $\alpha$ kills $\mathcal{R}$, so it descends
  uniquely to $\widetilde{f}:\mathcal{U}\otimes\mathcal{V}\to\mathcal{W}$.}
  \label{fig:tensor-quotient-construction-universal}
\end{figure}

\begin{corollary} \label{cor:tensor-product-unique-up-to-isomorphism} \hlt{Tensor Products Are Unique up to Isomorphism}\\
If $(\mathcal{T},t)$ and $(\mathcal{S},s)$ are tensor products of $\mathcal{U}$ and $\mathcal{V}$,\\
then there is a unique isomorphism $\Phi:\mathcal{T}\to\mathcal{S}$ satisfying $\Phi(t(\mathbf{u},\mathbf{v}))=s(\mathbf{u},\mathbf{v})$, where $\mathbf{u}\in\mathcal{U},\ \mathbf{v}\in\mathcal{V}$.
\begin{proof}
Apply universal property of $(\mathcal{T},t)$ to bilinear map $s$ and universal property of $(\mathcal{S},s)$ to bilinear map $t$.\\
This gives linear maps $\Phi:\mathcal{T}\to\mathcal{S}$ and $\Psi:\mathcal{S}\to\mathcal{T}$. The composites $\Psi\Phi$ and $I_{\mathcal{T}}$ are both mediating maps for $t$, while $\Phi\Psi$ and $I_{\mathcal{S}}$ are both mediating maps for $s$. Uniqueness forces $\Psi\Phi=I_{\mathcal{T}}$ and $\Phi\Psi=I_{\mathcal{S}}$.\\
\hlt{Structural Commentary}: Universal objects are unique up to canonical isomorphism.\\
\hlt{Key Insight}: Two universal factorisations force the comparison maps to be inverse.
\end{proof}
\end{corollary}

\begin{theorem} \label{thm:tensor-bases-dimension-and-rank} \hlt{Tensor Bases, Dimension, and Coordinate Rank}\\
Let $\mathcal{U}$ and $\mathcal{V}$ be finite-dimensional with bases $\beta=(\mathbf{u}_1,\ldots,\mathbf{u}_m)$ and $\gamma=(\mathbf{v}_1,\ldots,\mathbf{v}_n)$.
\begin{enumerate}[label=(\roman*)]
\setlength{\itemsep}{0pt}
\item The tensors $\mathbf{u}_i\otimes\mathbf{v}_j$ form a basis of $\mathcal{U}\otimes\mathcal{V}$.
\item $\dim(\mathcal{U}\otimes\mathcal{V})=\dim(\mathcal{U})\dim(\mathcal{V})$.
\item Every $\mathbf{z}\in\mathcal{U}\otimes\mathcal{V}$ has a unique matrix $R_{\beta,\gamma}(\mathbf{z})=[r_{ij}]\in\operatorname{Mat}(m,n,\mathbb{F})$ such that $\mathbf{z}=\sum\limits_{i=1}^{m}\sum\limits_{j=1}^{n} r_{ij}(\mathbf{u}_i\otimes\mathbf{v}_j)$.\\
If $\beta',\gamma'$ are other bases, then $\operatorname{rank}R_{\beta',\gamma'}(\mathbf{z})=\operatorname{rank}R_{\beta,\gamma}(\mathbf{z})$.
\end{enumerate}
\begin{proof}
\begin{enumerate}[label=(\roman*)]
\setlength{\itemsep}{0pt}
\item The coordinate tensor model of Definition \ref{def:finite-coordinate-tensor-product-model} has the universal property in Definition \ref{def:tensor-product-universal-property}.\\
Hence Corollary \ref{cor:tensor-product-unique-up-to-isomorphism} identifies $\mathcal{U}\otimes\mathcal{V}\cong\operatorname{Mat}(m,n,\mathbb{F})$ with $\mathbf{u}_i\otimes\mathbf{v}_j\mapsto E_{i,j}$, and $\{E_{i,j}\}$ is a basis.
\item Counting the $mn$ basis tensors gives $\dim(\mathcal{U}\otimes\mathcal{V})=mn=\dim(\mathcal{U})\dim(\mathcal{V})$.
\item The basis in (i) gives unique coefficients $r_{ij}$, hence $R_{\beta,\gamma}(\mathbf{z})$. If $\beta',\gamma'$ are other bases,\\
then $R_{\beta',\gamma'}(\mathbf{z})=P R_{\beta,\gamma}(\mathbf{z})Q$ for invertible change-of-basis matrices $P,Q$, so $\operatorname{rank}R_{\beta',\gamma'}(\mathbf{z})=\operatorname{rank}R_{\beta,\gamma}(\mathbf{z})$.
\end{enumerate}
\hlt{Structural Commentary}: Tensor rank is coordinate-free matrix rank.\\
\hlt{Key Insight}: Changing bases performs row and column changes, which preserve matrix rank.
\end{proof}
\end{theorem}

\begin{definition} \label{def:tensor-coordinate-matrix-and-rank} \hlt{Tensor Coordinate Matrix and Rank}\\
Let $\mathcal{U},\mathcal{V}$ be finite-dimensional, with bases $\beta=(\mathbf{u}_1,\ldots,\mathbf{u}_m)$ and $\gamma=(\mathbf{v}_1,\ldots,\mathbf{v}_n)$.
\begin{enumerate}[label=(\roman*)]
\setlength{\itemsep}{0pt}
\item For $\mathbf{z}=\sum\limits_{i=1}^{m}\sum\limits_{j=1}^{n} r_{ij}(\mathbf{u}_i\otimes\mathbf{v}_j)$, define the \hlt{coordinate matrix} $R_{\beta,\gamma}(\mathbf{z})\defeq[r_{ij}]$.
\item The \hlt{rank of the tensor} $\mathbf{z}$ is $\operatorname{rank}(\mathbf{z})\defeq\operatorname{rank}R_{\beta,\gamma}(\mathbf{z})$, independent of $\beta,\gamma$ by Theorem \ref{thm:tensor-bases-dimension-and-rank}.
\end{enumerate}
\end{definition}

\begin{theorem} \label{thm:tensor-zero-test-and-decomposable-rank} \hlt{Zero Tests and Decomposable Rank}\\
Let $\mathcal{U}$ and $\mathcal{V}$ be vector spaces.
\begin{enumerate}[label=(\roman*)]
\setlength{\itemsep}{0pt}
\item $\forall \mathbf{u}\in\mathcal{U},\mathbf{v}\in\mathcal{V}$, $\mathbf{0}\otimes\mathbf{v}=\mathbf{u}\otimes\mathbf{0}=\mathbf{0}$.
\item If $\mathbf{u}_1,\ldots,\mathbf{u}_r$ are linearly independent and $\sum\limits_{i=1}^{r}\mathbf{u}_i\otimes\mathbf{v}_i=\mathbf{0}$, then $\mathbf{v}_1=\cdots=\mathbf{v}_r=\mathbf{0}$.
\item If $\dim\mathcal{U},\dim\mathcal{V}<\infty$, then $\mathbf{z}\neq\mathbf{0}$ is decomposable $\Longleftrightarrow \operatorname{rank}(\mathbf{z})=1$.\\
More generally, $\operatorname{rank}(\mathbf{z})=\min\left\{r\ge0:\mathbf{z}=\sum\limits_{i=1}^{r}\mathbf{x}_i\otimes\mathbf{y}_i,\ \mathbf{x}_i\in\mathcal{U},\mathbf{y}_i\in\mathcal{V}\right\}$.
\end{enumerate}
\begin{proof}
\begin{enumerate}[label=(\roman*)]
\setlength{\itemsep}{0pt}
\item The tensor map $t:\mathcal{U}\times\mathcal{V}\to\mathcal{U}\otimes\mathcal{V}$ is bilinear, hence linear in each coordinate; therefore $t(\mathbf{0},\mathbf{v})=t(\mathbf{u},\mathbf{0})=\mathbf{0}$.
\item For each $k$, extend the dual basis functional on $\operatorname{span}(\mathbf{u}_1,\ldots,\mathbf{u}_r)$ to $\varphi_k\in\mathcal{U}'$ with $\varphi_k(\mathbf{u}_i)=\delta_{ki}$.\\
The bilinear map $(\mathbf{u},\mathbf{v})\mapsto\varphi_k(\mathbf{u})\mathbf{v}$ induces $T_k\in\LL(\mathcal{U}\otimes\mathcal{V},\mathcal{V})$.\\
Applying $T_k$ gives $\mathbf{0}=T_k\left(\sum_{i=1}^{r}\mathbf{u}_i\otimes\mathbf{v}_i\right)=\sum\limits_{i=1}^{r}\delta_{ki}\mathbf{v}_i=\mathbf{v}_k$.
\item Choose bases $\beta,\gamma$. If $\mathbf{z}=\mathbf{x}\otimes\mathbf{y}\neq\mathbf{0}$, then $R_{\beta,\gamma}(\mathbf{z})=[\mathbf{x}]_\beta[\mathbf{y}]_\gamma^{\mathsf{T}}$, so $\operatorname{rank}(\mathbf{z})=1$.\\
Conversely, if $\operatorname{rank}R_{\beta,\gamma}(\mathbf{z})=r$, write $R_{\beta,\gamma}(\mathbf{z})=\sum\limits_{i=1}^{r}a_i b_i^{\mathsf{T}}$ with columns $a_i,b_i$.\\
Translating $a_i,b_i$ back to vectors $\mathbf{x}_i,\mathbf{y}_i$ gives $\mathbf{z}=\sum\limits_{i=1}^{r}\mathbf{x}_i\otimes\mathbf{y}_i$ as a sum of $r$ decomposable tensors.\\
Conversely, any sum of $s$ decomposable tensors gives a sum of $s$ rank-one coordinate matrices, so $\operatorname{rank}(\mathbf{z})\le s$; hence the displayed minimum equals $\operatorname{rank}(\mathbf{z})$.
\end{enumerate}
\hlt{Structural Commentary}: Linear functionals can test one tensor coordinate at a time.\\
\hlt{Key Insight}: Tensor rank is matrix rank after choosing bases, but the number itself survives changing bases.
\end{proof}
\end{theorem}

\begin{theorem} \label{thm:tensor-linearizes-bilinear-maps} \hlt{Bilinearity Equals Linearity on a Tensor Product}\\
For $\mathbb{F}$-vector spaces $\mathcal{U},\mathcal{V},\mathcal{W}$, the map $\Phi:\operatorname{Bil}_{\mathbb{F}}(\mathcal{U},\mathcal{V};\mathcal{W})\to\LL(\mathcal{U}\otimes\mathcal{V},\mathcal{W})$, $\Phi(f)=\widetilde{f}$,\\
where $\widetilde{f}(\mathbf{u}\otimes\mathbf{v})=f(\mathbf{u},\mathbf{v})$, is an isomorphism.
\begin{proof}
The universal property gives existence and uniqueness of $\widetilde{f}$.\\
Define $\Psi:\LL(\mathcal{U}\otimes\mathcal{V},\mathcal{W})\to\operatorname{Bil}_{\mathbb{F}}(\mathcal{U},\mathcal{V};\mathcal{W})$ by $\Psi(L)\defeq L\circ t$.\\
Then $\Psi(\Phi(f))=\widetilde{f}\circ t=f$, and $\Phi(\Psi(L))=L$ by uniqueness of the mediating map for $L\circ t$. Hence $\Phi$ is bijective.\\
For $a\in\mathbb{F}$ and $f,g$ bilinear, $\Phi(af+g)(\mathbf{u}\otimes\mathbf{v})=af(\mathbf{u},\mathbf{v})+g(\mathbf{u},\mathbf{v})=(a\Phi(f)+\Phi(g))(\mathbf{u}\otimes\mathbf{v})$, so $\Phi$ is linear.\\
\hlt{Structural Commentary}: Tensor products package bilinear maps as ordinary linear maps.\\
\hlt{Key Insight}: To define a linear map out of $\mathcal{U}\otimes\mathcal{V}$, define a bilinear rule on $\mathcal{U}\times\mathcal{V}$.
\end{proof}
\end{theorem}

\begin{method} \label{meth:defining-linear-maps-on-tensor-products} \hlt{Defining Maps on Tensor Products}\\
To define $L\in\LL(\mathcal{U}\otimes\mathcal{V},\mathcal{W})$:
\begin{enumerate}[label=(\roman*)]
\setlength{\itemsep}{0pt}
\item define a candidate rule $f:\mathcal{U}\times\mathcal{V}\to\mathcal{W}$ by $f(\mathbf{u},\mathbf{v})$;
\item prove $f\in\operatorname{Bil}_{\mathbb{F}}(\mathcal{U},\mathcal{V};\mathcal{W})$;
\item set $L\defeq\Phi(f)$ from Theorem \ref{thm:tensor-linearizes-bilinear-maps}, so $L(\mathbf{u}\otimes\mathbf{v})=f(\mathbf{u},\mathbf{v})$.
\end{enumerate}
Do not assign arbitrary values on decomposable tensors: they span $\mathcal{U}\otimes\mathcal{V}$ but are usually linearly dependent.
\end{method}

\begin{theorem} \label{thm:tensor-product-of-linear-maps} \hlt{Tensor Product of Linear Maps}\\
Let $A\in\LL(\mathcal{U},\mathcal{U}')$ and $B\in\LL(\mathcal{V},\mathcal{V}')$.\\
There exists a unique $A\otimes B:\mathcal{U}\otimes\mathcal{V}\to\mathcal{U}'\otimes\mathcal{V}'$ such that $(A\otimes B)(\mathbf{u}\otimes\mathbf{v})=A\mathbf{u}\otimes B\mathbf{v}$.\\
Moreover, $(A_2A_1)\otimes(B_2B_1)=(A_2\otimes B_2)(A_1\otimes B_1)$, $I_{\mathcal{U}}\otimes I_{\mathcal{V}}=I_{\mathcal{U}\otimes\mathcal{V}}$.\\
In finite dimensions, $[A\otimes B]_{\beta\otimes\gamma}^{\beta'\otimes\gamma'}=[A]_{\beta}^{\beta'}\otimes[B]_{\gamma}^{\gamma'}$ for lexicographically ordered tensor-product bases, where the right-hand $\otimes$ is the matrix Kronecker product.
\begin{proof}
The map $(\mathbf{u},\mathbf{v})\mapsto A\mathbf{u}\otimes B\mathbf{v}$ is bilinear, so Theorem \ref{thm:tensor-linearizes-bilinear-maps} gives the unique $A\otimes B$.\\
The two functorial identities hold on decomposable tensors $((A_2\otimes B_2)(A_1\otimes B_1))(\mathbf{u}\otimes\mathbf{v})=A_2A_1\mathbf{u}\otimes B_2B_1\mathbf{v}$ and $(I_{\mathcal{U}}\otimes I_{\mathcal{V}})(\mathbf{u}\otimes\mathbf{v})=\mathbf{u}\otimes\mathbf{v}$, and decomposable tensors span. If $A\mathbf{u}_j=\sum_i a_{ij}\mathbf{u}'_i$ and $B\mathbf{v}_\ell=\sum_k b_{k\ell}\mathbf{v}'_k$, then $(A\otimes B)(\mathbf{u}_j\otimes\mathbf{v}_{\ell}) =\sum_{i,k}a_{ij}b_{k\ell}(\mathbf{u}'_i\otimes\mathbf{v}'_k)$, which is the Kronecker block-entry rule.\\
\hlt{Structural Commentary}: Tensoring maps acts independently in the two tensor coordinates.\\
\hlt{Key Insight}: The mixed product property for Kronecker products is the coordinate shadow of functoriality.
\end{proof}
\end{theorem}

\begin{theorem} \label{thm:tensor-product-of-map-spaces} \hlt{Tensor Product of Spaces of Linear Maps}\\
Let $\mathcal{U},\mathcal{V},\mathcal{X},\mathcal{Y}$ be $\mathbb{F}$-vector spaces. The assignment $(A,B)\mapsto A\otimes B$ induces a natural injective map\\
$\Theta:\LL(\mathcal{U},\mathcal{X})\otimes\LL(\mathcal{V},\mathcal{Y}) \to \LL(\mathcal{U}\otimes\mathcal{V},\mathcal{X}\otimes\mathcal{Y})$ with $\Theta(A\otimes B)(\mathbf{u}\otimes\mathbf{v})=A\mathbf{u}\otimes B\mathbf{v}$.\\
If $\mathcal{U},\mathcal{V},\mathcal{X},\mathcal{Y}$ are finite-dimensional, then $\Theta$ is an isomorphism.\\
In particular, $\mathcal{U}'\otimes\mathcal{V}'\hookrightarrow(\mathcal{U}\otimes\mathcal{V})'$; $(\varphi\otimes\psi)(\mathbf{u}\otimes\mathbf{v})=\varphi(\mathbf{u})\psi(\mathbf{v})$ is an isomorphism when $\dim\mathcal{U},\dim\mathcal{V}<\infty$.
\begin{proof}
\begin{enumerate}[label=(\roman*)]
\setlength{\itemsep}{0pt}
\item By Theorem \ref{thm:tensor-product-of-linear-maps}, $(A,B)\mapsto A\otimes B$ is bilinear, hence induces $\Theta$.
\item Let $0\neq Z=\sum\limits_{i=1}^{r}A_i\otimes B_i$ with $r$ minimal. Then $A_i\neq0$ and $(B_i)$ is linearly independent.\\
Choose $\mathbf{u}$ with some $A_i\mathbf{u}\neq0$, after reindexing let $A_1\mathbf{u},\ldots,A_m\mathbf{u}$ be basis of $\operatorname{span}\{A_i\mathbf{u}:1\le i\le r\}$.\\
Absorb remaining coefficients into $B_i$ gives $C_1,\ldots,C_m$, not all zero, s.t. $\Theta(Z)(\mathbf{u}\otimes\mathbf{v})=\sum\limits_{i=1}^{m}(A_i\mathbf{u})\otimes C_i\mathbf{v}$.\\
If $\Theta(Z)=0$, Theorem \ref{thm:tensor-zero-test-and-decomposable-rank} gives $C_i\mathbf{v}=0$ for all $i,\mathbf{v}$, contradiction. Thus $\ker\Theta=0$.
\item If $\mathcal{U},\mathcal{V},\mathcal{X},\mathcal{Y}$ are finite-dimensional, then the domain and codomain of $\Theta$ have equal dimension by Theorem \ref{thm:tensor-bases-dimension-and-rank}, so injectivity implies bijectivity. Taking $\mathcal{X}=\mathcal{Y}=\mathbb{F}$, $\LL(\mathcal{U},\mathbb{F})=\mathcal{U}'$, $\LL(\mathcal{U}\otimes\mathcal{V},\mathbb{F})=(\mathcal{U}\otimes\mathcal{V})'$ gives the dual embedding and finite-dimensional isomorphism.
\end{enumerate}
\hlt{Structural Commentary}: Tensoring maps yields maps on tensor products.\\
\hlt{Key Insight}: The rule on decomposable tensors is enough, and finite dimensions make the embedding onto.
\end{proof}
\end{theorem}

\begin{proposition} \label{prop:scalar-extension-by-tensor-product} \hlt{Scalar Extension by Tensor Product}\\
Let $K/\mathbb{F}$ be a field extension and let $\mathcal{V}_{\mathbb{F}}$ be an $\mathbb{F}$-vector space.\\
Then $K\otimes_{\mathbb{F}}\mathcal{V}_{\mathbb{F}}$ is a $K$-vector space via $\alpha\cdot(\beta\otimes\mathbf{v})\defeq(\alpha\beta)\otimes\mathbf{v}$. With $\mu:\mathcal{V}_{\mathbb{F}}\to K\otimes_{\mathbb{F}}\mathcal{V}_{\mathbb{F}}$, $\mu(\mathbf{v})\defeq1\otimes\mathbf{v}$, every $\mathbb{F}$-linear $f:\mathcal{V}_{\mathbb{F}}\to\mathcal{Y}$ into a $K$-vector space $\mathcal{Y}$ admits a unique $K$-linear $\widetilde{f}:K\otimes_{\mathbb{F}}\mathcal{V}_{\mathbb{F}}\to\mathcal{Y}$ such that $\widetilde{f}\circ\mu=f$.
\begin{proof}
For $\alpha\in K$, $(\beta,\mathbf{v})\mapsto(\alpha\beta)\otimes\mathbf{v}$ is $\mathbb{F}$-bilinear, hence induces $M_\alpha\in\LL_{\mathbb{F}}(K\otimes_{\mathbb{F}}\mathcal{V}_{\mathbb{F}})$.\\
On decomposable tensors, $M_\alpha M_\lambda=M_{\alpha\lambda}$, $M_1=I$, and distributivity holds; since decomposable tensors span, these define a $K$-vector-space structure.\\
For $f$, set $g(\beta,\mathbf{v})\defeq\beta f(\mathbf{v})$. Then $g$ is $\mathbb{F}$-bilinear, so it induces $\widetilde{f}$ with $\widetilde{f}(\beta\otimes\mathbf{v})=\beta f(\mathbf{v})$. This is $K$-linear and gives $\widetilde{f}\circ\mu=f$. If $h$ is $K$-linear and $h\circ\mu=f$, then $h(\beta\otimes\mathbf{v})=\beta h(1\otimes\mathbf{v})=\beta f(\mathbf{v})$, so $h=\widetilde{f}$.\\
\hlt{Structural Commentary}: Scalar extension tensors in the larger field.\\
\hlt{Key Insight}: Tensoring with $K$ upgrades which scalars are allowed to act.
\end{proof}
\end{proposition}
